:\chapter{Marco Teórico}
\thispagestyle{fancy}
\section{Definición de términos clave}
\noindent En este apartado se presentarán conceptos y términos clave que permitirán comprender y entender de mejor manera el contenido del documento.
\begin{itemize}

\item ChatBot: Para este protyecto se utiliza el concepto de chatbot, el cual podemos definir como, un sistema con el que el usuario puede interactuar generando una conversacion en lenguaje natural, haciendo preguntas acerca de diversos temas por medio de voz o texto, Este sistema podra aprender, entender y adaptarse a un contexto de diversos temas, ademas el chatbot podra hacer recomendaciones acerca de informacion que el usuario necesite o le pueda ser de ayuda\cite{chatbot_AWS,chatbot_IBM}.

\item Apuntes: En este proyecto haremos uso del concepto de apuntes, entendiendolo como una herramienta individual de aprendizaje personal, que promueve el aprendizaje al hacer uso de anotaciones breves y personales sobre clases, conferencias o lecturas para asi registrar ideas clave o resumir puntos importantes y nos dan una sintesis para un repaso efectivo\cite{Apuntes}. 

\item Notas: El uso del concepto de Nota podemos entenderlo como anotaciones complejas y organizadas que nos son utiles para la comprension de un tema estas pueden ayudar a complementar o mejorar los apuntes, presenta un mayor desarrollo en la informacion\cite{nota}. 

\item Frontend: Parte de una aplicacion con la que el usuario interactuara directamente, ademas se puede entender como todo lo que el usuario ve como: pantallas, botones, tablas, menus, etc\cite{frontend}.  

\end{itemize}

\section{Antecedentes}
\noindent En este apartado se presentan algunos antecedentes relacionados con el desarrollo de aplicaciones de gestión de notas académicas y sistemas de organización de información.   
\begin{itemize}
    \item \textbf{Evernote:} Es una aplicación popular para la toma y organización de notas que permite a los usuarios crear, editar y buscar notas de manera eficiente. Evernote ofrece funcionalidades como la categorización automática de notas y la capacidad de compartir notas entre usuarios.\cite{evernote,espana_evernote_nodate}

    \item \textbf{OneNote:} Desarrollada por Microsoft, OneNote es otra aplicación ampliamente utilizada para la gestión de notas académicas. Ofrece una interfaz intuitiva y fácil de usar, facilitando la creación y organización de notas, así como la colaboración entre estudiantes y profesores.\cite{one_note}

    \item \textbf{Notion:} Notion es una herramienta de productividad que combina la gestión de notas con otras funcionalidades como bases de datos y tareas. Es utilizada tanto por estudiantes como por profesionales para organizar información de manera eficiente.\cite{Notion_2025,Notion_Personal}

    \item \textbf{Google Keep:} Es una aplicación de Google que permite a los usuarios tomar notas rápidas y organizarlas mediante etiquetas y colores. Google Keep se integra con otros servicios de Google, facilitando el acceso a las notas desde diferentes dispositivos.\cite{google_keep}
    
    \item \textbf{Obsidian:} Aplicacion gratuita y flexible para organizar conocimiento, desde llevar la crea-
    cion de diarios, hasta crear bases de conocimientos y gestionar proyectos, obsidian crea los archivos
    en formatos abiertos por lo que no se depende de un solo formato ademas estos archivos los guarda
    de manera privada.
    

\end{itemize}   

Estos antecedentes demuestran la importancia y la demanda de aplicaciones de gestión de notas académicas, así como la necesidad de contar con herramientas que faciliten la organización y el acceso a la información de manera eficiente.
\section{Bases teóricas}
\noindent En este apartado se presentan las bases teóricas que sustentan el desarrollo de la aplicación de gestión de notas académicas.
\begin{itemize}
    \item \textbf{Gestión de Notas Académicas:} La gestión de notas académicas implica la organización, almacenamiento y acceso eficiente a la información relacionada con el aprendizaje. Esto incluye la creación de notas, la categorización de información y la capacidad de buscar y recuperar notas de manera rápida y efectiva.\cite{williams_reference_2024}  
    \item \textbf{Sistemas de Organización de Información:} Los sistemas de organización de información son herramientas que permiten a los usuarios estructurar y categorizar datos de manera eficiente. Estos sistemas pueden incluir funcionalidades como etiquetas, carpetas y categorías para facilitar la búsqueda y recuperación de información.\cite{collado_concepto_nodate}  
    \item \textbf{Interfaz de Usuario Intuitiva:} Una interfaz de usuario intuitiva es fundamental para garantizar que los usuarios puedan interactuar con la aplicación de manera efectiva. Esto implica el diseño de una interfaz clara y fácil de usar, que facilite la navegación y la realización de tareas relacionadas con la gestión de notas.\cite{ingenieria_de_usabilidad}  
\end{itemize}
Estas bases teóricas proporcionan el marco conceptual necesario para el desarrollo de una aplicación de gestión de notas académicas que sea eficiente, fácil de usar y que satisfaga las necesidades de los usuarios.
\vspace{1cm}


\section{Procesamiento de Lenguaje Natural}
El Procesamiento de Lenguaje Natural (PLN) es una rama de la inteligencia artificial que se enfoca en la interacción entre las computadoras y el lenguaje humano. El objetivo del PLN es permitir que las máquinas comprendan, interpreten y generen lenguaje natural de manera que sea útil para los usuarios. En el contexto de la gestión de notas académicas, el PLN puede ser utilizado para mejorar la organización y búsqueda de información mediante técnicas como el análisis de texto, la clasificación automática y la generación de resúmenes.\cite{procesamiento_lenguaje_natural}

\section{Modelos de Lenguaje}
Los modelos de lenguaje son algoritmos que utilizan técnicas de aprendizaje automático para predecir y generar
texto basado en patrones aprendidos a partir de grandes conjuntos de datos. Estos modelos pueden ser utilizados para diversas tareas, como la traducción automática, la generación de texto y la respuesta a preguntas. En el contexto de la gestión de notas académicas, los modelos de lenguaje pueden ser empleados para mejorar la capacidad de búsqueda y recuperación de información, así como para generar resúmenes automáticos de notas.\cite{llm} 
\section{Bases de Datos NoSQL}
Las bases de datos NoSQL son sistemas de gestión de bases de datos que no utilizan el modelo relacional tradicional. Estas bases de datos están diseñadas para manejar grandes volúmenes de datos no estructurados y ofrecen una mayor flexibilidad en términos de escalabilidad y rendimiento. En el contexto de la gestión de notas académicas, las bases de datos NoSQL pueden ser utilizadas para almacenar y organizar notas de manera eficiente, permitiendo un acceso rápido y flexible a la información.\cite{basededatosnosql}

\vspace{1cm}
\section{Sistemas de Recomendación}
Los sistemas de recomendación son herramientas que utilizan algoritmos para sugerir contenido relevante a los usuarios basándose en sus preferencias y comportamientos anteriores. Estos sistemas pueden ser utilizados en aplicaciones de gestión de notas académicas para recomendar recursos de estudio, artículos o notas relacionadas que puedan ser de interés para el usuario. Al implementar sistemas de recomendación, se puede mejorar la experiencia del usuario y facilitar el acceso a información relevante.\cite{sistemas_recomendacion}    

\vspace{1cm}
\section{Tecnologías Web}
Las tecnologías web son fundamentales para el desarrollo de aplicaciones modernas, incluyendo aquellas destinadas a la gestión de notas académicas. Estas tecnologías incluyen lenguajes de programación como HTML, CSS y JavaScript, así como frameworks y bibliotecas que facilitan la creación de interfaces de usuario interactivas y responsivas. Al utilizar tecnologías web, se puede garantizar que la aplicación de gestión de notas sea accesible desde diferentes dispositivos y plataformas, mejorando la experiencia del usuario.\cite{teconologias_web}   

\vspace{1cm}
\section{Metodologías Ágiles}
Las metodologías ágiles son enfoques de desarrollo de software que se centran en la colaboración, la flexibilidad y la entrega continua de valor al cliente. Estas metodologías promueven la iteración rápida, la retroalimentación constante y la adaptación a los cambios en los requisitos del proyecto. Al adoptar metodologías ágiles en el desarrollo de la aplicación de gestión de notas académicas, se puede garantizar que el proyecto se mantenga alineado con las necesidades de los usuarios y se entregue de manera eficiente.\cite{user_metodologias_2024}

\vspace{1cm}
\section{Seguridad y Privacidad de Datos}
La seguridad y privacidad de los datos son aspectos críticos en el desarrollo de aplicaciones que manejan información personal, como es el caso de la gestión de notas académicas. Es fundamental implementar medidas de seguridad para proteger los datos de los usuarios contra accesos no autorizados, pérdidas o robos. Además  
es importante garantizar la privacidad de los datos, cumpliendo con las regulaciones y normativas aplicables. Al abordar estos aspectos, se puede generar confianza en los usuarios y asegurar la integridad de la información manejada por la aplicación.\cite{noauthor_privacidad_nodate}  

\vspace{1cm}
\section{Experiencia de Usuario (UX)}
La experiencia de usuario (UX) se refiere a la percepción y respuesta de los usuarios al interactuar con una aplicación o sistema. Una buena UX es esencial para garantizar que los usuarios puedan utilizar la aplicación de manera efectiva y satisfactoria. En el contexto de la gestión de notas académicas, es importante diseñar una interfaz de usuario intuitiva, fácil de navegar y que facilite la realización de tareas relacionadas con la gestión de notas. Al centrarse en la UX, se puede mejorar la adopción y el uso de la aplicación por parte de los usuarios.\cite{ux}  
\vspace{1cm}
\section{Accesibilidad Digital}
La accesibilidad digital se refiere al diseño y desarrollo de aplicaciones y sitios web que sean utilizables por personas con discapacidades. Esto incluye la implementación de características que permitan a los usuarios con discapacidades visuales, auditivas, motoras o cognitivas acceder y utilizar la aplicación de manera efectiva. En el contexto de la gestión de notas académicas, es importante considerar la accesibilidad digital para garantizar que todos los usuarios, independientemente de sus capacidades, puedan beneficiarse de la aplicación.\cite{porto_accesibilidad_2023}

\vspace{1cm}
\section{Colaboración en Línea}
La colaboración en línea es un aspecto clave en el desarrollo de aplicaciones modernas, especialmente en el ámbito académico. Permitir que los usuarios colaboren en la creación, edición y organización de notas puede mejorar la eficiencia y la efectividad del proceso de aprendizaje. Al implementar funcionalidades de colaboración en línea en la aplicación de gestión de notas académicas, se puede fomentar el trabajo en equipo y facilitar la compartición de información entre estudiantes y profesores.\cite{trabajo_colaborativo}    

\vspace{1cm}
\section{Almacenamiento en la Nube}
El almacenamiento en la nube es una tecnología que permite a los usuarios guardar y acceder a sus datos desde cualquier lugar y en cualquier momento, utilizando servicios en línea. En el contexto de la gestión de notas académicas, el almacenamiento en la nube puede ofrecer ventajas significativas, como la accesibilidad remota, la escalabilidad y la seguridad de los datos. Al utilizar servicios de almacenamiento en la nube para la aplicación de gestión de notas, se puede mejorar la experiencia del usuario y garantizar la disponibilidad de la información en todo momento.\cite{almacenamiento_nube}

\vspace{1cm}
\section{Inteligencia Artificial en la Educación}
La inteligencia artificial (IA) está transformando el ámbito educativo al ofrecer soluciones innovadoras para mejorar el aprendizaje y la gestión de la información. En el contexto de la gestión de notas académicas, la IA puede ser utilizada para automatizar tareas como la categorización de notas, la generación de resúmenes y la personalización de recomendaciones de estudio. Al integrar tecnologías de IA en la aplicación de gestión de notas, se puede ofrecer una experiencia más eficiente y adaptada a las necesidades individuales de los usuarios.\cite{ia_educacion}
\vspace{1cm}
\section{Análisis de Sentimiento}
El análisis de sentimiento es una técnica de procesamiento de lenguaje natural que se utiliza para identificar y extraer opiniones y emociones expresadas en el texto. En el contexto de la gestión de notas académicas, el análisis de sentimiento puede ser utilizado para evaluar la retroalimentación de los usuarios sobre la aplicación, identificando áreas de mejora y comprendiendo mejor las necesidades de los usuarios. Al implementar análisis de sentimiento, se puede mejorar continuamente la aplicación y adaptarla a las expectativas de los usuarios.\cite{analisis_sentimiento}

% \section{}

